% GENERATED FILE. Edit canonical lessons, then run npm run generate:books.
% canonical-source-hash: fnv1a64:0b54f5bf87d206f7
% canonical-lessons: BN-C95-idaning, BN-C95-saradin, BN-C95-sararat, BN-C95-dupurbela, BN-C95-sokalbela

\chapter{Lately, All day, All night, Midday hours, Morning hours}
\label{ch:bn-lately-all-day-all-night-midday-hours-morning-hours}
\hlchaptermodality{95}

\begin{chapteropening}
\textbf{By the end of this chapter:} \emph{I can say lately, all day, all night, the midday hours and the morning hours.}

Complete the last lesson of chapter 95: I can say lately, all day, all night, the midday hours and the morning hours.
\end{chapteropening}

\section[\texorpdfstring{\bn{ইদানীং} (idānīng)}{ইদানীং (idānīng)}]{\bn{ইদানীং} (idānīng) — lately}

\label{lesson:BN-C95-idaning}



\begin{quote}
Before the new one: say the Bengali for ever, then the Bengali for soon.
\end{quote}

\subsection*{You'll want to know: \bn{ইদানীং}}
\textbf{\bn{ইদানীং}} — \emph{idānīng} — \textquotedblleft{}lately\textquotedblright{}.

\begin{cousinweb}
\textbf{\bn{ইদানীং}} (\emph{idānīng}) is Sanskrit \emph{idānīm}, \textquotedblleft{}now, at this time\textquotedblright{}, taken over whole as a learned loan, a tatsama.
\end{cousinweb}

\begin{grammarlens}[title={What you've built}]
One more word for this chapter.
\end{grammarlens}

\subsection*{Your turn}
\begin{itemize}
\raggedright
  \item \emph{Say it:} \emph{idānīng}
  \item \emph{Say it:} \emph{idānīng}, once more
  \item \emph{Say it:} say \emph{idānīng}
  \item \emph{Recall:} say the Bengali for ever, then the Bengali for soon, then say \emph{idānīng} again
\end{itemize}

\begin{tcolorbox}[breakable,colback=teal!4,colframe=teal!35!black,title={Before you move on}]
What does \bn{ইদানীং} mean? (\textquotedblleft{}Lately\textquotedblright{}.) Say it once more.
\end{tcolorbox}
\section[\texorpdfstring{\bn{সারাদিন} (sārādin)}{সারাদিন (sārādin)}]{\bn{সারাদিন} (sārādin) — all day}

\label{lesson:BN-C95-saradin}



\begin{quote}
Before the new one: say the Bengali for soon, then the Bengali for lately.
\end{quote}

\subsection*{You'll want to know: \bn{সারাদিন}}
\textbf{\bn{সারাদিন}} — \emph{sārādin} — \textquotedblleft{}all day\textquotedblright{}.

\begin{cousinweb}
\textbf{\bn{সারাদিন}} (\emph{sārādin}) is \emph{sārā}, \textquotedblleft{}whole\textquotedblright{}, + \emph{din}, \textquotedblleft{}day\textquotedblright{}. CDIAL lists Bengali \emph{sārā}, \textquotedblleft{}whole\textquotedblright{}, under Sanskrit \emph{sāra}, \textquotedblleft{}hardness, best part\textquotedblright{} (CDIAL 13355), and \emph{din} is Sanskrit \emph{dina}, \textquotedblleft{}day\textquotedblright{} (CDIAL 6328).
\end{cousinweb}

\begin{grammarlens}[title={What you've built}]
One more word for this chapter.
\end{grammarlens}

\subsection*{Your turn}
\begin{itemize}
\raggedright
  \item \emph{Say it:} \emph{sārādin}
  \item \emph{Say it:} \emph{sārādin}, once more
  \item \emph{Say it:} say \emph{sārādin}
  \item \emph{Recall:} say the Bengali for soon, then the Bengali for lately, then say \emph{sārādin} again
\end{itemize}

\begin{tcolorbox}[breakable,colback=teal!4,colframe=teal!35!black,title={Before you move on}]
What does \bn{সারাদিন} mean? (\textquotedblleft{}All day\textquotedblright{}.) Say it once more.
\end{tcolorbox}
\section[\texorpdfstring{\bn{সারারাত} (sārārāt)}{সারারাত (sārārāt)}]{\bn{সারারাত} (sārārāt) — all night}

\label{lesson:BN-C95-sararat}



\begin{quote}
Before the new one: say the Bengali for lately, then the Bengali for all day.
\end{quote}

\subsection*{You'll want to know: \bn{সারারাত}}
\textbf{\bn{সারারাত}} — \emph{sārārāt} — \textquotedblleft{}all night\textquotedblright{}.

\begin{cousinweb}
\textbf{\bn{সারারাত}} (\emph{sārārāt}) is \emph{sārā}, \textquotedblleft{}whole\textquotedblright{}, + \emph{rāt}, \textquotedblleft{}night\textquotedblright{}. CDIAL lists Bengali \emph{sārā}, \textquotedblleft{}whole\textquotedblright{}, under Sanskrit \emph{sāra}, \textquotedblleft{}hardness, best part\textquotedblright{} (CDIAL 13355), and \emph{rāt} under \emph{rātrī}, \textquotedblleft{}night\textquotedblright{} (CDIAL 10702).
\end{cousinweb}

\begin{grammarlens}[title={What you've built}]
One more word for this chapter.
\end{grammarlens}

\subsection*{Your turn}
\begin{itemize}
\raggedright
  \item \emph{Say it:} \emph{sārārāt}
  \item \emph{Say it:} \emph{sārārāt}, once more
  \item \emph{Say it:} say \emph{sārārāt}
  \item \emph{Recall:} say the Bengali for lately, then the Bengali for all day, then say \emph{sārārāt} again
\end{itemize}

\begin{tcolorbox}[breakable,colback=teal!4,colframe=teal!35!black,title={Before you move on}]
What does \bn{সারারাত} mean? (\textquotedblleft{}All night\textquotedblright{}.) Say it once more.
\end{tcolorbox}
\section[\texorpdfstring{\bn{দুপুরবেলা} (dupurbelā)}{দুপুরবেলা (dupurbelā)}]{\bn{দুপুরবেলা} (dupurbelā) — the midday hours}

\label{lesson:BN-C95-dupurbela}



\begin{quote}
Before the new one: say the Bengali for all day, then the Bengali for all night.
\end{quote}

\subsection*{You'll want to know: \bn{দুপুরবেলা}}
\textbf{\bn{দুপুরবেলা}} — \emph{dupurbelā} — \textquotedblleft{}the midday hours\textquotedblright{}.

\begin{cousinweb}
\textbf{\bn{দুপুরবেলা}} (\emph{dupurbelā}) is \emph{dupur}, \textquotedblleft{}midday\textquotedblright{}, usually explained as Sanskrit \emph{dviprahara}, \textquotedblleft{}the second watch\textquotedblright{} (CDIAL 8900), + \emph{belā}, \textquotedblleft{}time\textquotedblright{}, from Sanskrit \emph{velā}, \textquotedblleft{}limit, time\textquotedblright{}, a word CDIAL thinks probably Dravidian (CDIAL 12115).
\end{cousinweb}

\begin{grammarlens}[title={What you've built}]
One more word for this chapter.
\end{grammarlens}

\subsection*{Your turn}
\begin{itemize}
\raggedright
  \item \emph{Say it:} \emph{dupurbelā}
  \item \emph{Say it:} \emph{dupurbelā}, once more
  \item \emph{Say it:} say \emph{dupurbelā}
  \item \emph{Recall:} say the Bengali for all day, then the Bengali for all night, then say \emph{dupurbelā} again
\end{itemize}

\begin{tcolorbox}[breakable,colback=teal!4,colframe=teal!35!black,title={Before you move on}]
What does \bn{দুপুরবেলা} mean? (\textquotedblleft{}The midday hours\textquotedblright{}.) Say it once more.
\end{tcolorbox}
\section[\texorpdfstring{\bn{সকালবেলা} (sôkālbelā)}{সকালবেলা (sôkālbelā)}]{\bn{সকালবেলা} (sôkālbelā) — the morning hours}

\label{lesson:BN-C95-sokalbela}



\begin{quote}
Before the new one: say the Bengali for all night, then the Bengali for the midday hours.
\end{quote}

\subsection*{You'll want to know: \bn{সকালবেলা}}
\textbf{\bn{সকালবেলা}} — \emph{sôkālbelā} — \textquotedblleft{}the morning hours\textquotedblright{}.

\begin{cousinweb}
\textbf{\bn{সকালবেলা}} (\emph{sôkālbelā}) is \emph{sôkāl}, \textquotedblleft{}morning\textquotedblright{}, which matches Sanskrit \emph{sakāla}, \textquotedblleft{}seasonable\textquotedblright{} (CDIAL 13067), + \emph{belā}, \textquotedblleft{}time\textquotedblright{}, from Sanskrit \emph{velā}, \textquotedblleft{}limit, time\textquotedblright{}, a word CDIAL thinks probably Dravidian (CDIAL 12115).
\end{cousinweb}

\begin{grammarlens}[title={What you've built}]
One more word for this chapter.
\end{grammarlens}

\subsection*{Your turn}
\begin{itemize}
\raggedright
  \item \emph{Say it:} \emph{sôkālbelā}
  \item \emph{Say it:} \emph{sôkālbelā}, once more
  \item \emph{Say it:} say \emph{sôkālbelā}
  \item \emph{Recall:} say the Bengali for all night, then the Bengali for the midday hours, then say \emph{sôkālbelā} again
\end{itemize}

\begin{tcolorbox}[breakable,colback=teal!4,colframe=teal!35!black,title={Before you move on}]
What does \bn{সকালবেলা} mean? (\textquotedblleft{}The morning hours\textquotedblright{}.) Say it once more.
\end{tcolorbox}
